\input{ieee-preamble}
\hypersetup{
  pdftitle={Transfer and Shared Learning for Atari Control},
  pdfauthor={Santiago Restrepo},
  pdfsubject={Pixel-based control, feature transfer, and balanced shared reinforcement learning},
  pdfkeywords={Atari, reinforcement learning, transfer learning, multi-task learning, continual learning}
}

\title{Transfer and Shared Learning for Atari Control}
\author{\IEEEauthorblockN{Santiago Restrepo}
\IEEEauthorblockA{\href{mailto:contact@waajacu.com}{contact@waajacu.com}\quad
\url{https://waajacu.com/}\\[3pt]
\textit{Shared perception, specialized control, and measurable retention}}}
\date{}

\begin{document}
\maketitle

\begin{abstract}
This approach learns Atari control from pixels through policy
optimization, feature transfer, and balanced practice across games.
A shared visual encoder supports specialized policies. The design
separates perception from control, isolates the effect of transferred
knowledge, and checks retention explicitly. Native parallel simulation,
independent playback, and preserved experiment evidence make the
learning process efficient and inspectable.
\end{abstract}

\begin{IEEEkeywords}
Atari, reinforcement learning, transfer learning, shared representation.
\end{IEEEkeywords}

\balance
\section{Learn Control from Visual Experience}
The Arcade Learning Environment (ALE) supplies screens, actions, and
game rewards~\cite{machado2018}. Grayscale conversion, resizing, and
temporal pooling produce compact observations. Stacking recent screens
exposes motion to a convolutional encoder without object coordinates
or handwritten gameplay rules. A registry records each game's ROM
identity, action count, and evaluation criterion.

The shared encoder feeds a game-specific actor and critic:
\begin{equation}
z_t=f_\theta(o_t),\qquad
a_t\sim\pi_{\phi_g}(\,\cdot\mid z_t),\qquad
v_t=V_{\psi_g}(z_t).
\end{equation}
Here $o_t$ is the screen stack, $g$ selects the game's heads, and $v_t$
estimates future discounted training reward. Shared visual features
remain separate from game-specific action meanings.

\section{Improve Policies with Game Rewards}
Proximal Policy Optimization (PPO) alternates on-policy interaction
with minibatch learning~\cite{schulman2017}. Actions are sampled from
the actor. Generalized advantage estimation uses rewards and critic
predictions to identify better-than-expected actions. A clipped policy
objective discourages excessive changes, critic regression improves
prediction, and an entropy incentive encourages exploration.

Rewards accumulate over repeated actions and are clipped for learning;
evaluation retains raw returns. Success thresholds are evaluation
criteria, not reward bonuses. Life loss preserves episode continuity.
Time-limit truncation bootstraps from the critic rather than being
treated as a natural terminal outcome.

\section{Isolate the Effect of Feature Transfer}
Initialize a target encoder from a learned source policy. Create fresh
target actor and critic heads and a fresh optimizer, then fine-tune
the complete network. This reuses perception without importing
incompatible actions or source optimizer history.

Compare against learning from scratch with matched target interaction
budgets, seeds, initial heads, and evaluation conditions. Account for
source pretraining separately. This isolates whether visual reuse
helps adaptation instead of assuming that transferred weights help.

\section{Balanced Shared Learning and Retention}
To add a game, restore the shared encoder and existing actor--critic
heads, initialize new heads, and start a fresh joint optimizer.
Collect fresh experience from earlier games alongside the new game,
keeping their objectives active throughout adaptation.

Balance environments and minibatches across games, normalize
advantages within each game, and weight game losses equally. This
limits domination by interaction volume or advantage scale while
allowing every task to update shared features. Compare each policy
before and after adaptation: continued practice supports retention,
but conflicting gradients can still cause forgetting.

\section{Efficient and Inspectable Execution}
Native C++ workers advance emulators on the CPU; batched inference and
optimization use the GPU through LibTorch. Independent browser
playback loads saved policies without pacing training. A reproducible
Debian environment separates dependency installation from lifecycle
and experiment commands.

Preserve joint checkpoints, game-specific inference exports, settings,
and source lineage. Atomic checkpoint publication lets concurrent
readers access complete artifacts without coupling playback to
the learner's mutable state.

\section{Make Evaluation a Design Contract}
Freeze policies and use consistent preprocessing, sticky actions,
reset rules, and full-game termination. Evaluation selects the most
probable action. Pair seeds for retention comparisons and reserve
separate seeds for final confirmation. Distinguish natural endings
from capped episodes; assess each game's criterion independently.
Preserve hashes and episode traces so each claim identifies its
policy, training lineage, and evaluation protocol.

\begingroup
\raggedright
\bibliographystyle{IEEEtran}
\bibliography{references}
\endgroup
\end{document}
